\chapter{Zusammenfassung und Ausblick}
\label{ch:07_zusammenfassung_ausblick}

% Einleitung.
Im Folgenden werden die in der vorliegenden Ausarbeitung gewonnenen Erkenntnisse in einem abschließenden Fazit unter
Rückbesinnung auf die Kern-Forschungsfragen zusammengefasst.
Es erfolgt die Ableitung eines Resümees bezüglich der Leistungsfähigkeit der untersuchten Rendering-Paradigmen im
Kontext der selektiven Stilisierung.
Auf dieser Grundlage werden konkrete Anforderungen für zukünftige Forschungs- und Entwicklungsarbeiten abgeleitet,
insbesondere in Bezug auf das entwickelte Softwareartefakt (\ac{HyDra}-Framework).


\section{Fazit der Arbeit}
\label{sec:0701_conclusion}

% Prämisse.
In der vorliegenden Untersuchung wurde der fundamentale, architektonische Konflikt zwischen der Skalierbarkeit komplexer
Echtzeit-Rendering-Szenen unter hoher Stil-Entropie für selektives Non-Photorealistic Rendering (\ac{NPR}) adressiert.
Das eigens entwickelte Evaluationsframework (\ac{HyDra}) wurde als isolierte Testumgebung zur Analyse der
untersuchten Basis-Rendering-Paradigmen verwendet.
Die Zusammenfassungen orientieren sich an den in Kapitel ~\ref{ch:06_diskussion} dargelegten Ergebnissen.

% Vorgehen.
Das methodische Vorgehen basierte auf der Durchführung systematischer, empirischer Tests der Rendering-Paradigmen unter
kontrollierten Bedingungen.
Im Rahmen der Untersuchung wurden die geometrische Dichte, die Beleuchtungskomplexität, Auflösungen, Bandbreiten und
die räumliche Stil-Entropie quantifiziert.

% Messergebnisse.
Die empirischen Messungen identifizierten spezifische Hardware-Bottlenecks für jede der untersuchten
Rendering-Pipelines.
Die Single-Pass-Forward-Pipeline wies vor allem bei ansteigender Lichtkomplexität, durch \ac{ALU}-Last bedingt,
Schwächen auf.
Die Single-Pass-Forward-Pipeline wies vor allem bei ansteigender Lichtkomplexität aufgrund hoher \ac{ALU}-Last
deutliche Schwächen auf.
Die Single-Pass-Deferred-Uber-Pipeline zeigte bei hoher räumlicher Stil-Durchmischung signifikante Leistungseinbrüche,
die durch Warp-Divergenz und Registerdruck verursacht wurden.
Die Deferred-Volume-Pipeline unterlag bei stark überlappenden Lichtquellen einer Bandbreitenlimitierung durch erhöhte
\ac{RMW}-Zyklen.

% Architekturergebnisse.
Auf architektonischer Ebene zeigte die Single-Pass-Forward-Pipeline zudem die signifikante Einschränkung,
Bildraumeffekte nicht selektiv darstellen zu können.
Die Single-Pass-Deferred-Uber-Pipeline wies zudem Schwächen in der Lichtberechnung auf, was zu Pre-Light Zwängen in
der Verarbeitung bestimmter Kernel-basierter \ac{NPR}-Effekte geführt hat.
Insgesamt erwiesen sich Single-Pass Architekturen als unzureichend, um extrem hohe Lichtmengen darzustellen.

% Beantwortung der Forschungsfrage.
Die Synthese dieser Ergebnisse lässt die zentrale Forschungsfrage dahingehend beantworten, dass im Kontext der
selektiven Stilisierung keine universell überlegene Basis-Rendering-Architektur existiert.
Die Leistungsfähigkeiten wurde dynamisch und ausschließlich anhand der dominierenden Szenenlast determiniert.
Den zentralen praktischen Beitrag dieser Arbeit bildet das entwickelte \ac{HyDra}-Framework, welches Entwicklern eine
gekapselte, von proprietärem Engine-Overhead freie Testumgebung zur empirischen Evaluation von Rendering-Architekturen
bereitstellt.
Als abgeleitetes Resultat dieser Testumgebung wurde in Kapitel~\ref{sec:0602_heuristic_model} ein heuristisches
Entscheidungsmodell formuliert, das anhand deterministischer \textit{Hard-Limits} und empirischer \textit{Soft-Limits}
als konkretes Werkzeug zur frühzeitigen Architekturbewertung dient.


\section{Ausblick und zukünftige Arbeiten}
\label{sec:0702_future_work}

% Profiler.
Die nachfolgenden Ansätze bauen größtenteils auf den in Kapitel~\ref{sec:0603_reflektion_and_limits} identifizierten
Limitationen auf.
Aus den methodischen Schwächen der reinen Zeitmessung folgt beispielsweise für Folgestudien die Notwendigkeit einer
erweiterten Telemetrie-Erfassung.
Sollen die genauen Auswirkungen von Hardware-Phänomenen wie Registerdruck, Cache-Misses oder Speichersättigung
evaluiert werden, ist die Integration systemnaher Hardware-Profiler wie \textit{NVIDIA Nsight} oder der \textit{Radeon
Developer Tool Suite (RDTS)} erforderlich.

% Pipeline-Erweiterung.
Im Rahmen der architektonischen Weiterentwicklung des \ac{HyDra}-Frameworks sind Ergänzungen weiterer
Rendering-Pipelines von zentralem Interesse.
Neben einer simplen Multi-Pass-Forward-Pipeline sind hier etwa
Forward+~\parencite[vgl.][]{haradaForwardBringingDeferred2012},
Clustered-Forward/Deferred~\parencite[vgl.][]{olssonClusteredDeferredForward2012},
Tiled~Forward/Deferred~\parencite[vgl.][]{olssonTiledClusteredForward2012} oder Deferred-Pipelines mit
Visibility-Buffern~\parencite[vgl.][]{burnsVisibilityBufferCacheFriendly2013} mögliche relevante weitere Pipelines.

% Vektor-Erweiterung.
In diesem Kontext sind insbesondere die Implementierung unterschiedlicher Lichtquellen (Richtungs- und Flächenlichter)
und das damit einhergehende Shadow-Mapping sowie erweiterte Materialeigenschaften (Spiegelungen und Transluzenzen) von
primärer Relevanz.

% Effekt-Untersuchungen.
Zudem lassen sich neben der stilistischen Entropie durch komplexere Materialeigenschaften auch die Einflüsse von
Anti-Aliasing-Verfahren im Kontext der untersuchten Rendering-Paradigmen analysieren.
Dies betrifft insbesondere die selektive Stilisierung, da die verursachten Artefakte wie harte Materialkanten das
Erscheinungsbild maßgeblich beeinflussen.
Anti-Aliasing kann somit als weitere semantische Evaluationsgrundlage der Rendering-Pipelines herangezogen werden.

% Hardware.
Abschließend ist auch eine Portierung der Evaluationssoftware auf unterschiedliche Systeme und Hardwarebedingungen von
Interesse, da der Kern der Evaluationen in Form von Hardwarekipppunkten unter anderen Hardwarearchitekturen grundlegend
anders ausfallen kann.
Dadurch wird es möglich, spezifische Hardware-Phänomene präziser zu isolieren und zu untersuchen.

